\documentclass[11pt]{article}
\usepackage[margin=1in]{geometry}
\usepackage{amsmath,amssymb,amsthm,mathtools}
\usepackage[T1]{fontenc}
\usepackage{lmodern,microtype}
\usepackage[hidelinks]{hyperref}
\hypersetup{pdftitle={An explicit counterexample to three-summand Kato-kernel splitting},
pdfauthor={Alec Kriebel},pdfsubject={A negative answer to OWR 5/2016, Question 2},
pdfkeywords={Kato cohomology, Brauer kernels, characteristic two, Pfister forms}}
\newtheorem{theorem}{Theorem}
\newtheorem{proposition}[theorem]{Proposition}
\newtheorem{remark}[theorem]{Remark}
\newcommand{\F}{\mathbb F}
\newcommand{\K}{\mathcal K}
\newcommand{\dlog}{\operatorname{dlog}}
\newcommand{\Br}{\operatorname{Br}}
\newcommand{\N}{\operatorname{N}}
\title{An explicit counterexample to three-summand Kato-kernel splitting}
\author{Alec Kriebel\thanks{Independent researcher. ORCID:
\href{https://orcid.org/0009-0001-9320-500X}{0009-0001-9320-500X}.}}
\date{17 October 2026\\[2pt]\small Unrefereed, extensively AI-assisted research note}
\begin{document}
\maketitle
\begin{abstract}
Question 2 in Laghribi's contribution to Oberwolfach Report 5/2016 asks
whether a mixed Kato kernel is the sum of the kernels of a finite purely
inseparable extension, a separable quadratic extension, and a bilinear
Pfister function field. We give an explicit negative instance in
\(H_2^2\), using an exponent-two extension and a binary bilinear Pfister
form that remains anisotropic. The mixed class is exactly a specialization
of the generators of Aravire--Laghribi--O'Ryan (2019/2020). We prove directly
that it remains outside the sum after the additional Pfister kernel is
included, by a universal norm and residue obstruction. The original
degree-one cohomology case is separately positive. No firstness claim is made.
\end{abstract}

\section{The question and the counterexample}
For a field \(E\) of characteristic two, put
\[
 H_2^{n+1}(E)=
 \Omega_E^n/\bigl(\wp(\Omega_E^n)+d\Omega_E^{n-1}\bigr),\qquad
 H_2^1(E)=E/\wp(E).
\]
Here \(\wp\) on logarithmic generators is induced by
\(u\,d\log v_1\wedge\cdots\mapsto
(u^2+u)\,d\log v_1\wedge\cdots\), modulo exact forms.
Write \(\K_n(E/F)\) for the kernel of
\(H_2^{n+1}(F)\to H_2^{n+1}(E)\), and abbreviate
\(\K_1(E/F)=\K(E/F)\).

The original question \cite[printed p.~251, Question~2]{OWR} asks whether
\begin{equation}\label{question}
 \K_n(LK(B)/F)=\K_n(L/F)+\K_n(K/F)+\K_n(F(B)/F)
\end{equation}
for every finite purely inseparable \(L/F\), separable quadratic \(K/F\),
and anisotropic bilinear Pfister form \(B\) remaining anisotropic over \(L\).
As in that source, \(F(B)\) is the function field of the
\emph{projective quasilinear quadric} of \(v\mapsto B(v,v)\).
This convention includes an anisotropic binary form, whose quadric is
integral and zero-dimensional.

The positive exponent-one regime is covered by \cite{ALO16}.
Higher-exponent two-field nonsplitting and the mixed generators are
antecedent work \cite{ALO20}. The class below is exactly their
Definition~1.3 and Lemma~3.1 specialization: in their notation take
\(n=m=1\), \(f_0=1\), \(f_1=z\), and \(a=s^{-4}\).
The only nonzero generator term is
\((\beta^2z/(1+\beta^2z))^2d\log z\).
Moreover \(s^{-4}+s^{-1}=\wp(s^{-2}+s^{-1})\), so these parameters define
the same separable quadratic field. This note supplies the direct
three-summand nonmembership proof; it makes no priority claim for
the class, example, or target deduction.

\begin{theorem}\label{main}
Let \(b,z,s,t\) be algebraically independent over \(\F_2\), and set
\[
 F=\F_2(b,z,s,t),\quad L=F(\beta),\quad\beta^4=b,\qquad
 K=F(\theta),\quad\theta^2+\theta=s^{-1},\qquad
 B=\langle1,t\rangle_b .
\]
Then \(L/F\) has degree four and exponent two, \(K/F\) is separable
quadratic, and \(B_L\) is anisotropic. The class
\[
 \xi=\overline{s^{-4}\frac{bz^2}{1+bz^2}\,d\log z}\ \in H_2^2(F)
\]
satisfies
\[
 \xi\in\K(LK/F)\subseteq\K(LK(B)/F),\qquad
 \xi\notin\K(L/F)+\K(K/F)+\K(F(B)/F).
\]
Consequently \eqref{question} is false at its allowed index \(n=1\).
\end{theorem}

We use the natural identification \(H_2^2(E)\simeq\Br(E)[2]\),
sending \(u\,d\log v\) to the characteristic-two quaternion class
\([u,v)_E\); see \cite[pp.~2--3]{CQM}.
Its generators satisfy \(i^2+i=u,\ j^2=v,\ ji=(i+1)j\).
For a genuine Artin--Schreier quadratic field \(E(\eta)/E\),
\(\eta^2+\eta=u\), the cyclic relative Brauer theorem says
\[
 \K(E(\eta)/E)=\{[u,c)_E:c\in E^*\}.
\]
Indeed a normalized two-cocycle of the order-two Galois group is
specified by an invariant scalar \(c\); coboundaries multiply it by a
norm. Its crossed product is the displayed quaternion algebra
\cite[IV, \S3, items~3.11--3.14]{Milne}.
Also \([u,c)_E=0\) if and only if \(c\) is a norm from \(E(\eta)^*\)
\cite[Theorem~6.4.11]{Voight}.
These statements allow imperfect base fields.

\section{Proof of Theorem \ref{main}}
The \(b\)-valuation shows that \(b\) is nonsquare, so \(X^4-b\)
is irreducible and \(L=\F_2(\beta,z,s,t)\) has the claimed degree and
exponent. A negative valuation of \(x^2+x\) is even; hence the odd
\(s\)-pole of \(s^{-1}\) proves the Artin--Schreier extension is genuine.
The \(t\)-valuation in \(L\) proves \(t\notin L^2\), giving anisotropy.
It remains nonsquare in separable algebraic extensions of \(L\).

Let \(\gamma^2=t\). The quadric \(X^2+tY^2=0\) in \(\mathbb P^1\)
is entirely in the \(Y\ne0\) chart, with field coordinate ring
\(F[X/Y]/((X/Y)^2+t)\). Thus
\[
 F(B)=F(\gamma),\qquad
 M:=LF(B)=\F_2(\beta,z,s,\gamma),\qquad LK(B)=M(\theta).
\]
The same descriptions over \(K\) are valid since a separable
extension cannot contain a new purely inseparable square root.

\paragraph{Vanishing in the composite.}
In \(L\), put \(f=1+\beta^2z\) and \(h=\beta^2z/f\).
Since \(d(\beta^2)=0\),
\[
 d\log f=h\,d\log z,\qquad h^2=\frac{bz^2}{1+bz^2}.
\]
Apply the Artin--Schreier relation first to
\(s^{-2}h\,d\log z\), and then to \(s^{-1}d\log f\). In cohomology,
\begin{equation}\label{vanish}
 \xi_L=\overline{s^{-2}h\,d\log z}
       =[s^{-2},f)_L=[s^{-1},f)_L .
\end{equation}
This is zero over \(LK\). The equalities occur in the quotient by
exact forms and Artin--Schreier images, not by treating \(\wp\)
as \(L\)-linear.

\paragraph{Reduction of all three summands.}
Suppose \(\xi=x+y+w\), with the three individual kernel memberships
in Theorem~\ref{main}. The cyclic description gives
\(y=[s^{-1},c)_F\) for some \(c\in F^*\).
Restrict to \(M=LF(B)\): both \(x\) and \(w\) vanish.
By \eqref{vanish} and multiplicativity of the second symbol slot,
\begin{equation}\label{normcondition}
 [s^{-1},f/c)_M=0.
\end{equation}
Thus the extra Pfister summand is fully included in the argument.

\paragraph{A necessary residue condition for every \(c\).}
Put \(k=\F_2(b,z,t)\), \(k'=\F_2(\beta,z,\gamma)\).
Complete \(M=k'(s)\) to \(D=k'((s))\).
The odd-pole argument also proves that \(T=D(\theta)\) is a quadratic
field. With \(\pi=\theta^{-1}\),
\[
 \pi^2+s\pi+s=0,\qquad \N_{T/D}(\pi)=s,\qquad
 \N_{T/D}(A+B\pi)=A^2+sAB+sB^2 .
\]
Here the conjugate of \(\pi\) is \(\pi+s\).
For nonzero \(A,B\in D\), let \(a=v_s(A)\), \(b'=v_s(B)\).
The even integer \(2a\) and the odd integer \(1+2b'\) differ.
The middle-term valuation \(1+a+b'\) is strictly greater than their
minimum: split into \(a\le b'\) and \(a\ge b'+1\).
Hence the norm has valuation \(\min(2a,1+2b')\).
The cases \(A=0\) or \(B=0\) follow from the same formula.
It follows that a unit norm has \(v_s(A)=0,\ v_s(B)\ge0\),
and its residue is a square in \(k'\).
This proof uses no perfectness or defect assumption.

For arbitrary \(c\in F^*=k(s)^*\), write
\(c=s^r c_0(1+O(s))\) with \(r\in\mathbb Z,\ c_0\in k^*\).
Equation \eqref{normcondition} makes \(f/c\) a norm over \(D\).
Multiplication by the norm \(s^r=\N(\pi^r)\), valid also for negative
\(r\), makes it a unit norm. Its residue therefore satisfies
\begin{equation}\label{square}
 \frac{1+\beta^2z}{c_0}\in k'^{*2}.
\end{equation}

\paragraph{The residue condition is impossible.}
Set
\[
 S=k'^2=\F_2(\beta^2,z^2,\gamma^2),\qquad
 R=\F_2(\beta^4,z^2,\gamma^2).
\]
We have \(k=R(z)\), of degree two over \(R\), so
\(c_0=u+vz\) uniquely with \(u,v\in R\), including arbitrary rational
denominators. Also \(z\notin S\), as seen by the odd valuation of
\(z^2\) in the rational field \(S\). Comparing \(1,z\) coefficients in
\[
 1+\beta^2z=q(u+vz),\qquad q\in S^*,
\]
gives \(qu=1,\ qv=\beta^2\). Thus \(u\ne0\) and
\(\beta^2=v/u\in R\), a contradiction: as an independent variable,
\(\beta^2\) has valuation one at zero, whereas nonzero elements of
\(R=\F_2((\beta^2)^2,z^2,\gamma^2)\) have even valuation there.
This disproves \eqref{square} for every \(c\), proving the theorem.
\qed

\section{The boundary degree and scope}
\begin{proposition}
At \(n=0\), equality \eqref{question} holds for every set of its stated
data with \(B\) a positive-fold bilinear Pfister form.
\end{proposition}
\begin{proof}
An Artin--Schreier root cannot first appear in a purely inseparable
extension, since that root is separable over the base. Nor can it first
appear in a rational function field, whose algebraic constants are its
base field. An anisotropic quasilinear form of dimension \(d\ge2\) has
projective-quadric function field purely inseparable
quadratic over a rational field in \(d-2\) variables: use a projective
chart and solve for one squared coordinate. For \(d=2\), there are no
rational variables.

Separable extension preserves quasilinear anisotropy, by linear
disjointness with purely inseparable extensions. Thus the chart
description applies over \(KL\). Any Artin--Schreier root over \(F\)
in \(LK(B)\) descends first to \(KL\), then to \(K\), since \(KL/K\)
is purely inseparable. Hence the composite kernel is the \(K/F\)
kernel, and the other two kernels are zero.
\end{proof}
If one separately assigns \(F(B)=F\) to a zero-fold form, its
\(n=0\) case follows from the same purely inseparable injection.
No function field of an empty projective quadric is needed.

Two boundaries of the construction are useful.
If the Pfister slot is \(z\) rather than the independent \(t\),
adjoining its field kills \(dz\) and absorbs \(\xi\).
If \(b=u^2\) already in the base and \(\beta^2=u\), then the computation
\emph{over \(L\)} gives \(\xi_L=[s^{-1},1+uz)_L\).
Since \(1+uz\) is now in the base, this proves membership in
\(\K(L/F)+\K(K/F)\). It does not justify a single-symbol equality in
the base: there \(d(1+uz)=u\,dz+z\,du\).
This corrects an auxiliary shortcut in the earlier draft and leaves
the exponent-two example unchanged.

\paragraph{Verification, credit and review status.}
The supplement gives exact algebra diagnostics, including the
fixed-logarithmic-basis exact-differential correction, and records
the mathematical and bounded primary priority audits.
Finite checks supplement the universal proof; they are not a
cohomology decision procedure. The audits did not verify an earlier
recorded full three-summand answer in the bodies inspected.
Final journal-body comparison for \cite{ALO20} and the full bodies
of the adjacent nonnormal-extension paper
(DOI \href{https://doi.org/10.1080/00927872.2019.1686510}{10.1080/00927872.2019.1686510})
and 2022 transfer paper
(DOI \href{https://doi.org/10.1016/j.jpaa.2021.106930}{10.1016/j.jpaa.2021.106930})
remain access limits. No global openness, firstness, or causal-independence
certificate is asserted.
AI tools were used extensively in developing, writing, reproducing and
adversarially checking this work. Initial internal review families had
parent context available; their checks were not perfectly isolated.
This preprint is unrefereed, has not received external human peer review,
and is not formally proof-assistant verified.

\begin{thebibliography}{9}
\bibitem{OWR}
A.~Laghribi (joint with R.~Aravire and M.~O'Ryan),
Differential and Quadratic Forms in Characteristic Two,
in \emph{Oberwolfach Report 5/2016}, 248--251.
\href{https://doi.org/10.4171/OWR/2016/5}{doi:10.4171/OWR/2016/5}.

\bibitem{ALO16}
R.~Aravire, A.~Laghribi and M.~O'Ryan,
Graded Witt kernels of the compositum of multiquadratic extensions
with the function fields of Pfister forms,
\emph{J. Algebra} \textbf{449} (2016), 635--659.
\href{https://doi.org/10.1016/j.jalgebra.2015.11.022}{doi:10.1016/j.jalgebra.2015.11.022}.

\bibitem{ALO20}
R.~Aravire, A.~Laghribi and M.~O'Ryan,
Cohomological kernels of mixed extensions in characteristic 2,
\emph{J. Algebra} \textbf{542} (2020), 249--276.
\href{https://doi.org/10.1016/j.jalgebra.2019.09.012}{doi:10.1016/j.jalgebra.2019.09.012}.
Definition/lemma numbering refers to the author manuscript dated
15 August 2019,
\href{https://laghribi.perso.math.cnrs.fr/cohom-kernels.pdf}{author PDF}.

\bibitem{CQM}
A.~Chapman and A.~Qu\'eguiner-Mathieu,
Minimal quadratic forms for the function field of a conic in characteristic 2,
\href{https://arxiv.org/abs/2111.15251v3}{arXiv:2111.15251v3} (2022).

\bibitem{Milne}
J.~S.~Milne, \emph{Class Field Theory}, version 4.03 (2020), Chapter IV.
\href{https://www.jmilne.org/math/CourseNotes/CFT.pdf}{Author course notes}.

\bibitem{Voight}
J.~Voight, \emph{Quaternion Algebras}, Graduate Texts in Mathematics 288,
Springer, 2021, Chapter 6.
\href{https://doi.org/10.1007/978-3-030-56694-4_6}{doi:10.1007/978-3-030-56694-4\_6}.
\end{thebibliography}
\end{document}
